
\documentclass{article}
\usepackage{colortbl}
\usepackage{makecell}
\usepackage{multirow}
\usepackage{supertabular}

\begin{document}

\newcounter{utterance}

\centering \large Interaction Transcript for game `hot\_air\_balloon', experiment `air\_balloon\_survival\_en\_reasoning off\_easy', episode 2 with Qwen3{-}30B{-}A3B{-}t0.0.
\vspace{24pt}

{ \footnotesize  \setcounter{utterance}{1}
\setlength{\tabcolsep}{0pt}
\begin{supertabular}{c@{$\;$}|p{.15\linewidth}@{}p{.15\linewidth}p{.15\linewidth}p{.15\linewidth}p{.15\linewidth}p{.15\linewidth}}
    \# & $\;$A & \multicolumn{4}{c}{Game Master} & $\;\:$B\\
    \hline

    \theutterance \stepcounter{utterance}  
    & & \multicolumn{4}{p{0.6\linewidth}}{
        \cellcolor[rgb]{0.9,0.9,0.9}{
            \makecell[{{p{\linewidth}}}]{
                \texttt{\tiny{[P1$\langle$GM]}}
                \texttt{You are participating in a collaborative negotiation game.} \\
\\ 
\texttt{Together with another participant, you must agree on a single set of items that will be kept. Each of you has your own view of how much each item matters to you (importance). You do not know how the other participant values the items. Additionally, you are given the effort each item demands.} \\
\texttt{You may only agree on a set if the total effort of the selected items does not exceed a shared limit:} \\
\\ 
\texttt{LIMIT = 3218} \\
\\ 
\texttt{Here are the individual item effort values:} \\
\\ 
\texttt{Item effort = \{"A19": 919, "B91": 852, "A30": 299, "B11": 762, "C47": 161, "C75": 205, "C26": 380, "C73": 399, "A47": 534, "A15": 333, "A75": 100, "A39": 420, "A56": 354, "C24": 130, "B26": 589\}} \\
\\ 
\texttt{Here is your personal view on the importance of each item:} \\
\\ 
\texttt{Item importance values = \{"A19": 138, "B91": 583, "A30": 868, "B11": 822, "C47": 783, "C75": 65, "C26": 262, "C73": 121, "A47": 508, "A15": 780, "A75": 461, "A39": 484, "A56": 668, "C24": 389, "B26": 808\}} \\
\\ 
\texttt{Goal:} \\
\\ 
\texttt{Your goal is to negotiate a shared set of items that benefits you as much as possible (i.e., maximizes total importance to YOU), while staying within the LIMIT. You are not required to make a PROPOSAL in every message {-} you can simply negotiate as well. All tactics are allowed!} \\
\\ 
\texttt{Interaction Protocol:} \\
\\ 
\texttt{You may only use the following structured formats in a message:} \\
\\ 
\texttt{PROPOSAL: \{'A', 'B', 'C', …\}} \\
\texttt{Propose keeping exactly those items.} \\
\texttt{REFUSE: \{'A', 'B', 'C', …\}} \\
\texttt{Explicitly reject opponent's proposal.} \\
\texttt{ARGUMENT: \{'...'\}} \\
\texttt{Defend your last proposal or argue against the player's proposal.} \\
\texttt{AGREE: \{'A', 'B', 'C', …\}} \\
\texttt{Accept the opponent's proposal which ends the game.} \\
\\ 
\\ 
\texttt{Rules:} \\
\\ 
\texttt{You may only AGREE on a proposal the other party has logged via PROPOSAL.} \\
\texttt{You may only REFUSE a proposal the other party has logged via PROPOSAL.} \\
\texttt{Total effort of any PROPOSAL or AGREE set must be ≤ LIMIT.} \\
\texttt{Do NOT reveal your hidden importance scores.} \\
\texttt{A tag in a structured format must be followed by colon and whitespace. The argument must be a python set containing 0 or more strings.} \\
\texttt{So, it must be of the form TAG: \{...\}} \\
\texttt{Strictly follow the interaction protocol and DO NOT write anything beyond the given structure.} \\
\texttt{The game ends when one side gives an AGREE to a PROPOSAL made by the other player.} \\
\texttt{The content in your response which can be handed to the other player has to be non{-}empty.} \\
\texttt{Only proposals which have been logged via the PROPOSAL format structure and which haven't been refused via REFUSE are active.} \\
\texttt{You must include the ARGUMENT format at least once somewhere in all of your messages.} \\
            }
        }
    }
    & \\ \\

    \theutterance \stepcounter{utterance}  
    & \multicolumn{4}{p{0.6\linewidth}}{
        \cellcolor[rgb]{0.9,0.9,0.9}{
            \makecell[{{p{\linewidth}}}]{
                \texttt{\tiny{[P1$\rangle$GM]}}
                \texttt{<think>} \\
\texttt{Okay, let's see. I need to figure out a set of items that maximizes my importance while keeping the total effort under 3218. The other participant's importance values are unknown, so I can't directly negotiate based on their preferences. But since we need to agree on a set, I should aim for a combination that's likely acceptable to both, maybe by focusing on items with high importance to me but not too high effort.} \\
\\ 
\texttt{First, I'll list out the items with their effort and my importance. Let me check the effort and importance values again.} \\
\\ 
\texttt{Item effort:} \\
\texttt{A19: 919, B91:852, A30:299, B11:762, C47:161, C75:205, C26:380, C73:399, A47:534, A15:333, A75:100, A39:420, A56:354, C24:130, B26:589.} \\
\\ 
\texttt{My importance:} \\
\texttt{A19:138, B91:583, A30:868, B11:822, C47:783, C75:65, C26:262, C73:121, A47:508, A15:780, A75:461, A39:484, A56:668, C24:389, B26:808.} \\
\\ 
\texttt{I need to select items where the sum of their efforts is ≤3218. Let's think about which items have high importance to me. For example, A30 (importance 868) and B11 (822) are high. But their efforts are 299 and 762. Let me try to pick items with high importance per effort ratio.} \\
\\ 
\texttt{Alternatively, maybe start by selecting items with the highest importance and check if they fit. Let's see:} \\
\\ 
\texttt{B91: importance 583, effort 852. If I include B91, that's 852. Then B} \\
            }
        }
    }
    & & \\ \\

    \theutterance \stepcounter{utterance}  
    & & \multicolumn{4}{p{0.6\linewidth}}{
        \cellcolor[rgb]{0.9,0.9,0.9}{
            \makecell[{{p{\linewidth}}}]{
                \texttt{\tiny{[P1$\langle$GM]}}
                \texttt{Your response did not contain an argument. You must include the structured format ARGUMENT: \{...\} somewhere in your response. Try again} \\
            }
        }
    }
    & \\ \\

    \theutterance \stepcounter{utterance}  
    & \multicolumn{4}{p{0.6\linewidth}}{
        \cellcolor[rgb]{0.9,0.9,0.9}{
            \makecell[{{p{\linewidth}}}]{
                \texttt{\tiny{[P1$\rangle$GM]}}
                \texttt{<think>} \\
\texttt{Okay, the user pointed out that my previous response didn't include an ARGUMENT. I need to make sure that in my next message, I include the ARGUMENT: \{...\} structure. Let me think about how to structure this properly.} \\
\\ 
\texttt{First, I need to propose a set of items. Let me try to select items that have high importance to me but manageable effort. Let's see:} \\
\\ 
\texttt{Looking at my importance values, the top items are A30 (868), B11 (822), A15 (780), C47 (783), B26 (808), A56 (668), A39 (484), A75 (461), A47 (508), etc. But I need to check their efforts.} \\
\\ 
\texttt{Let me try to pick some of these. For example, A30 (299), B11 (762), A15 (333), C47 (161), B26 (589). Let's add their efforts: 299 + 762 = 1061, +333 = 1394, +161 = 1555, +589 = 2144. That's under 3218. Then maybe add others like A56 (354) → 2144 + 354 = 2498. Then A39 (420) → 2918. Then A47 (534) would take it to 3452, which is over. So maybe not. Alternatively, add C24 (130) → 2918 +130=3048. Then maybe A75 (100) → 3148. Then maybe C26 (380) would be 3528, too much. So maybe that set: A30, B11, A15, C47, B26, A56, A39, C24, A75. Total effort 3148. That's under 3218. The total importance would be 868 +822 +} \\
            }
        }
    }
    & & \\ \\

    \theutterance \stepcounter{utterance}  
    & & \multicolumn{4}{p{0.6\linewidth}}{
        \cellcolor[rgb]{0.9,0.9,0.9}{
            \makecell[{{p{\linewidth}}}]{
                \texttt{\tiny{[P1$\langle$GM]}}
                \texttt{Your response contained an untagged sequence or you used STRATEGIC REASONING more than once. You may only use the structured formats as explained in the initial message. They must all be of the form TAG: \{...\}} \\
            }
        }
    }
    & \\ \\

    \theutterance \stepcounter{utterance}  
    & \multicolumn{4}{p{0.6\linewidth}}{
        \cellcolor[rgb]{0.9,0.9,0.9}{
            \makecell[{{p{\linewidth}}}]{
                \texttt{\tiny{[P1$\rangle$GM]}}
                \texttt{<think>} \\
\texttt{Okay, the user is strict about the format. I need to make sure that my response only includes the structured tags: PROPOSAL, REFUSE, ARGUMENT, AGREE. No other text. Also, I must include an ARGUMENT once.} \\
\\ 
\texttt{Let me try again. First, I need to propose a set. Let me pick items that have high importance to me and check the total effort. Let's take A30 (299), B11 (762), A15 (333), C47 (161), B26 (589), A56 (354), A39 (420), C24 (130), A75 (100). Adding up the efforts: 299+762=1061, +333=1394, +161=1555, +589=2144, +354=2498, +420=2918, +130=3048, +100=3148. That's under 3218. The importance would be 868+822+780+783+808+668+484+389+461. Let me calculate that: 868+822=1690, +780=2470, +783=3253, +808=4061, +668=4729, +484=5213, +389=5602, +461=6063. That's a high total. So the proposal would be these items. Now, I need to format it as PROPOSAL: \{'A30', 'B11', 'A15', 'C47', 'B26', 'A56', 'A39', 'C24', 'A75'\}, then add an ARGUMENT explaining why this set is good. The argument should be a set, maybe pointing out the high importance and effort within limit. So the message would be:} \\
\\ 
\texttt{PROPOSAL: \{'A3} \\
            }
        }
    }
    & & \\ \\

    \theutterance \stepcounter{utterance}  
    & & & \multicolumn{2}{p{0.3\linewidth}}{
        \cellcolor[rgb]{0.9,0.9,0.9}{
            \makecell[{{p{\linewidth}}}]{
                \texttt{\tiny{[GM$|$GM]}}
                \texttt{abort} \\
            }
        }
    }
    & & \\ \\

    \theutterance \stepcounter{utterance}  
    & & & \multicolumn{2}{p{0.3\linewidth}}{
        \cellcolor[rgb]{0.9,0.9,0.9}{
            \makecell[{{p{\linewidth}}}]{
                \texttt{\tiny{[GM$|$GM]}}
                \texttt{end game} \\
            }
        }
    }
    & & \\ \\

\end{supertabular}
}

\end{document}
